Volba mezi cloudovým enterprise řešením a lokálním / on‑premise provozem se řídí třemi hlavními kritérii: citlivostí zpracovávaných dat, požadavky na integraci do existujícího ekosystému (SSO, LMS, katalogy uživatelů) a provozní/rozpočtovou dostupností. Před rozhodnutím proveďte krátké posouzení rizik a rozsahu — zmapujte scénáře použití, datové toky a potenciální expozice. Zahrňte do posouzení také otázky provozní resilience, zálohování a požadavky na obnovu po havárii. Pokud v řešení budou zpracovávána osobní nebo citlivá studentská data, je doporučené provést DPIA jako součást plánování pilotu; výsledky DPIA pomohou specifikovat technické i smluvní požadavky vůči poskytovateli a zohlednit právní rizika \cite{kb_malinka_chatgpt_ve_vyuce_dva_roky_pote_2024_pdf_90bc3aaf_page_12_1}. DPIA rovněž slouží jako podklad pro interní rozhodovací procesy a pro komunikaci s vedením školy či fakultou.

Praktická doporučení pro rozhodovací proces (VUT)
\begin{enumerate}
\item Klasifikujte data podle citlivosti a účelu zpracování; vytvořte jednoduchou matici klasifikace (veřejné, interní, citlivé, zvlášť chráněné) a pro scénáře se studentskými údaji naplánujte DPIA. Zahrňte výstupy DPIA do požadavků na bezpečnostní opatření a do smluvního vyjednávání s poskytovatelem \cite{kb_malinka_chatgpt_ve_vyuce_dva_roky_pote_2024_pdf_90bc3aaf_page_12_1}. Taková klasifikace umožní rychlé rozhodování o tom, které datové sady mohou být zpracovávány v cloudu a které musí zůstat lokálně.
\item Preferujte enterprise/licencované varianty pro výukové scénáře, zejména pokud chcete využít integrované školní funkce a centralizovanou správu uživatelů. Enterprise balíčky často nabízejí administrativní nástroje, SLA a integrace (příklady: nástroje Microsoft 365 a školní licence obsahující edukační služby), které usnadní nasazení a podporu \cite{kb_ai_101_tipu_pdf_04e4a115_page_15_2,kb_ai_101_tipu_pdf_04e4a115_page_36_1}. Při posuzování variant porovnejte nejen cenu licence, ale i dostupnost podpory, školení a možnosti rozšíření funkcí.
\item Požadujte jasnou smlouvu o zpracování (DPA) a SLA zahrnující incident response, retenci dat, bezpečnostní standardy a auditní možnosti; výsledky DPIA jsou užitečným podkladem při specifikaci těchto požadavků. Dbejte na definice odpovědnosti za sub‑procesory a na dostupnost logů pro forenzní šetření \cite{kb_ai_101_tipu_pdf_04e4a115_page_8_1,kb_malinka_chatgpt_ve_vyuce_dva_roky_pote_2024_pdf_90bc3aaf_page_12_1}. V dokumentech si vyjasněte postupy oznámení bezpečnostních incidentů a lhůty pro nápravu.
\item Ověřte technické možnosti integrace: SSO/autorizace, centrální správa přístupů a auditní logy — tyto prvky usnadňují provoz a zajišťují sledovatelnost chování systému; zkontrolujte také podporu šifrování v klidu i při přenosu a možnosti exportu dat (general background knowledge). Zjistěte, jak snadno lze z řešení migrovat data či uzavřít službu bez ztráty provozních záznamů a metrik.
\item Pro piloty vyžadujte minimální množství PII v indexovaných zdrojích a tam, kde je to možné, pseudonymizujte nebo agregujte data před nahráním do externích služeb. Uvažujte o testovacích datech a anonymizovaných sadách, abyste snížili riziko úniku citlivých informací \cite{kb_ai_101_tipu_pdf_04e4a115_page_8_1}. Doporučuje se také zavést pravidelné kontroly indexovaných zdrojů a proces pro rychlé odstranění citlivých záznamů.
\end{enumerate}

Konkrétní kroky spolupráce s IT a právním oddělením (postup)
\begin{itemize}
\item Krátký assessment: garant předmětu a IT společně vyplní formulář s typem dat, očekávanými uživateli, účelem zpracování a navrhovanými integracemi; pokud se zpracovávají studentská data, iniciujte DPIA a zapojte právní oddělení dříve než v pozdějších fázích projektu \cite{kb_malinka_chatgpt_ve_vyuce_dva_roky_pote_2024_pdf_90bc3aaf_page_12_1}. Tento krok by měl zahrnovat i odhad nákladů a zdrojů potřebných pro pilotní provoz.
\item Smluvní příprava: na základě výsledků DPIA specifikujte požadavky na DPA/SLA (retence, sub‑procesory, auditní práva, incident response) a v případě potřeby vyžádejte hosting/rezidenci dat v přijatelné jurisdikci. Zahrňte také metriky dostupnosti, sankce za porušení a postupy pro bezpečnostní incidenty (general background knowledge; doporučeno zahrnout do DPA). Jasně definované smluvní podmínky usnadní řešení sporů a zkrátí dobu reakce v případě incidentu.
\item Technická integrace: nastavte SSO a role‑based access, zapněte auditní logování a metriky využití; v pilotní fázi měřte chybovost, výkonnost a bezpečnostní incidenty pro rozhodnutí o dalším rozšíření. Dokumentujte integrační kroky, zálohovací politiky a plán obnovy po havárii. Plánujte také pravidelné revize konfigurace a testy obnovy dat.
\item Testování funkcí a pedagogické ověření: ověřte, že funkce určené pro výuku (např. automatizovaná tvorba testů nebo specializované edukační moduly) fungují v rámci vybraného enterprise balíčku. Proveďte uživatelské testy s garanty předmětů a sbírejte zpětnou vazbu; konkrétním příkladem je Microsoft 365 Copilot ve spojení s Forms, který umí generovat testy a vysvětlení, což ukazuje integraci pedagogických funkcí do enterprise prostředí \cite{kb_ai_101_tipu_pdf_04e4a115_page_36_1}. Zpětnou vazbu dokumentujte a zapracujte do rozhodnutí o nasazení či úpravě integračních nastavení.
\end{itemize}

Poznámka k RAG a lokálnímu indexování
\begin{itemize}
\item Retrieval‑Augmented Generation (RAG) a indexace interních dokumentů mohou snížit potřebu ukládání PII přímo do externích modelů, ale vyžadují přísné řízení přístupů, revizi indexovaných zdrojů a kontrolu, jaké dokumenty jsou indexovány. Implementujte pravidla pro redakci citlivých polí, dočasné cache a minimalizaci kontextu zasílaného do modelu; při návrhu RAG architektury spolupracujte úzce s oddělením IT a právem, aby byly pokryty požadavky vyplývající z DPIA a DPA (general background knowledge). Zvažte rovněž možnosti auditu indexů a pravidelné automatizované kontroly kvality indexovaných dat.
\end{itemize}

Shrnutí pro garanty předmětů: při rozhodování upřednostněte enterprise/licencované řešení tam, kde chcete garantovat školní funkce, centrální správu uživatelů a SLA; tam, kde je kritická kontrola nad daty nebo jsou vyšší právní rizika, zvažte lokální/on‑premise variantu. V obou případech považujte DPIA za výchozí krok plánování pilotu a použijte její výstupy pro smluvní i technické požadavky vůči poskytovateli \cite{kb_malinka_chatgpt_ve_vyuce_dva_roky_pote_2024_pdf_90bc3aaf_page_12_1, kb_ai_101_tipu_pdf_04e4a115_page_15_2}. Dodržení těchto kroků zlepší bezpečnostní postavení projektu, zkrátí dobu nasazení a sníží provozní rizika spojená s nasazením AI nástrojů ve výuce.